\documentclass[../Project-Report-Main.tex]{subfiles}

\begin{document}

\chapter{Introduction}

\section{Motivation}

\begin{itemize}
 \item In recent years, digital financial transactions have experienced explosive growth, particularly in India driven by the Unified Payments Interface (UPI), mobile banking applications, and digital wallets. Millions of daily micro-transactions—ranging from grocery purchases and food delivery to utility bill payments—are processed digitally. While digital payments offer unprecedented convenience, they simultaneously create a significant challenge for personal financial management (PFM). Individuals frequently find themselves making scores of transactions monthly across multiple platforms (e.g., SBI YONO, Google Pay, PhonePe, Paytm), leading to fragmented records, lack of visibility into cumulative spending habits, and financial disorganization.

 \item Traditional expense tracking methods, such as manual ledger logging or static spreadsheet entry (e.g., Microsoft Excel, Google Sheets), require continuous discipline and manual effort. Research indicates that a majority of users abandon manual tracking apps within a few weeks due to user fatigue and administrative burden. Conversely, existing automated mobile apps often rely on reading private user SMS messages—raising severe privacy and security concerns—or fail to support native, password-protected Indian bank PDF statements.

 \item Furthermore, traditional PFM tools merely present static tables or basic charts without offering actionable, context-aware financial guidance. Users are left with raw data but lack intelligent advice on how to optimize their budget, reduce discretionary overspending, or improve net savings.

 \item This project, \textbf{Finlytics} (FinTracker), is motivated by the critical need for a privacy-focused, intelligent, full-stack personal finance tracker tailored to the modern financial ecosystem. By automating the parsing of password-protected Indian bank statements (such as SBI YONO) and multi-line Google Pay UPI PDFs, implementing a dynamic hybrid categorization engine, providing intuitive interactive visual analytics, and incorporating Large Language Model (LLM) powered financial advice, Finlytics bridges the gap between raw transaction data and meaningful, actionable financial decision-making.
\end{itemize}

\section{Problem Statement}

Modern personal finance tracking is hindered by several core challenges:

\begin{enumerate}
 \item \textbf{Complex \& Encrypted Statement Formats:} Indian bank PDF statements (e.g., SBI YONO) and UPI transaction logs (e.g., Google Pay) are structured heterogeneously. They often employ PDF password encryption (e.g., combination of user name and date of birth), multi-line transaction narrations, embedded bank reference codes, and varying tabular layouts. Standard off-the-shelf text extractors fail to parse these documents accurately.

 \item \textbf{High Friction of Manual Expense Tracking:} Manual entry of daily expenses leads to user fatigue, data omissions, human calculation errors, and eventual abandonment of tracking altogether.

 \item \textbf{Cryptic Transaction Narratives \& Poor Categorization:} Raw bank narrations contain extraneous technical codes and unorganized merchant details. Without intelligent parsing and customizable categorization rules, transactions remain uncategorized or incorrectly grouped, rendering expense analysis ineffective.

 \item \textbf{Static Visualizations without Actionable Insights:} Existing PFM solutions offer basic graphical summaries (bar charts or pie charts) but lack intelligent advisory capabilities. They do not interpret financial trends, identify top overspending channels, or offer personalized monetary advice tailored to the user's specific financial behavior.

 \item \textbf{Privacy Risks in Automated Tracking:} Commercial financial apps frequently demand intrusive permissions, such as full read-access to personal SMS messages or direct linking to bank accounts via open-banking APIs, exposing sensitive financial and identity data to third-party privacy breaches.
\end{enumerate}

\section{Proposed System}

The proposed system, \textbf{Finlytics}, is a robust, full-stack, web-based Personal Finance Management application built using Spring Boot, Spring Data JPA, Hibernate, MySQL, HTML5, CSS3, Vanilla JavaScript, Chart.js, Apache PDFBox, and the Groq API (LLaMA 3.1 8B Instant model).

Finlytics solves the aforementioned problems through the following integrated subsystems:

\begin{enumerate}
 \item \textbf{Automated Multi-Format Statement Parser:} A dedicated Java parsing module (\texttt{SbiStatementParser}) utilizing Apache PDFBox capable of loading password-protected bank PDFs (e.g., SBI YONO) and multi-line Google Pay UPI PDFs. The parser uses advanced Regular Expressions (Regex) to extract transaction dates, transaction types (WDL, DEP, CSH, TRF, UPI), debit/credit amounts, running balances, and clean merchant/narration strings while discarding header/footer noise.

 \item \textbf{Hybrid Rule-Based \& Learnt Categorization Engine:} A multi-tiered categorization mechanism that evaluates transaction descriptions against:
 \begin{itemize}
   \item User-defined persistent merchant rules stored in MySQL (\texttt{MerchantRule} entity).
   \item A fallback keyword dictionary mapping common merchant indicators (e.g., Swiggy/Zomato to "Food", Petrol/IRCTC/Uber to "Travel", Amazon/Flipkart to "Shopping").
   \item Default category assignments ("Salary" for income, "Personal" for expense) to ensure 100\% of transactions are classified.
 \end{itemize}

 \item \textbf{Dynamic Budgeting \& Visual Analytics Dashboard:} An interactive web interface powered by Chart.js featuring:
 \begin{itemize}
   \item Summary cards displaying Total Income, Total Expenses, and Net Savings.
   \item Categorical expense breakdown pie/doughnut charts and income vs. expense comparisons.
   \item Monthly trend graphs and top merchant spending leaderboards.
   \item Category budget tracking with dynamic visual progress bars and overspending warning flags.
 \end{itemize}

 \item \textbf{AI-Driven Financial Advisor (LLM Integration):} An intelligent financial advisory endpoint (\texttt{InsightsController}) that aggregates the user's last 30-day spending data (category breakdown, net savings, top merchants) into a structured prompt payload. It leverages the Groq API executing the \texttt{llama-3.1-8b-instant} model to generate concise, tailored, actionable monetary recommendations. Responses are cached per session to minimize API latencies and resource usage.

 \item \textbf{Privacy-Centric \& Secure Architecture:} Runs entirely under user control using session-based authentication (Register, Login, Logout) and isolated database tables, ensuring sensitive personal financial data remains private and secure without external SMS reading or open-banking data sharing.
\end{enumerate}

\section{Objectives}

The primary objectives of the Finlytics project are:

\begin{enumerate}
 \item To design and implement an automated bank statement parsing engine capable of handling password-protected PDF files (SBI YONO) and multi-line UPI statements (Google Pay).
 \item To build a robust transaction management system offering automated ingestion, manual CRUD operations, date-range filtering, and standard CSV data export.
 \item To develop a hybrid transaction categorization system incorporating persistent merchant-to-category rules and multi-keyword heuristic dictionaries.
 \item To implement a budget tracking module that allows users to define monthly spending limits per category and receive visual overspending alerts.
 \item To construct an intuitive, responsive web dashboard with interactive Chart.js analytics for clear financial visualization.
 \item To integrate Large Language Model (LLM) capabilities via Groq API (LLaMA 3.1) to generate automated, context-aware, personalized financial insights and actionable advice based on 30-day transactional trends.
 \item To enforce secure user authentication and strict data privacy through session-based authorization and user-isolated database storage.
\end{enumerate}

\pagebreak
\section{Methodology}

The development of Finlytics follows a structured software engineering lifecycle utilizing the Model-View-Controller (MVC) design pattern and RESTful web service architecture. The methodology comprises five key phases:

\begin{itemize}
 \item \textbf{Phase 1: Architecture \& Data Modeling}
 \begin{itemize}
   \item Relational database tables designed in MySQL mapped via Spring Data JPA entities: \texttt{User}, \texttt{Transaction}, \texttt{Category}, \texttt{Budget}, \texttt{MerchantRule}, and \texttt{Upload}.
 \end{itemize}

 \item \textbf{Phase 2: PDF Parser Development (Data Ingestion)}
 \begin{itemize}
   \item Apache PDFBox Integration: Uploaded files are ingested via Spring \texttt{MultipartFile} and unlocked using user-provided password parameters.
   \item Regex Pipeline Execution: Text tokenization via \texttt{PDFTextStripper}, date pattern matching (\texttt{DATE\_LINE}, \texttt{DATE\_INLINE\_LINE}), multi-line narrative assembly, amount parsing (\texttt{AMOUNT\_LINE}), and UPI merchant extraction (\texttt{UPI\_NAME}).
 \end{itemize}

 \item \textbf{Phase 3: Categorization \& Persistence Engine}
 \begin{itemize}
   \item Evaluation Order: Check \texttt{MerchantRule} repository $\rightarrow$ Iterate static keyword dictionaries $\rightarrow$ Assign system fallback category.
   \item Persistence: Store structured \texttt{Transaction} objects linked to the logged-in \texttt{User} entity in MySQL.
 \end{itemize}

 \item \textbf{Phase 4: Analytics \& AI Insights Endpoint}
 \begin{itemize}
   \item Aggregation Pipeline: Calculates 30-day rolling aggregate statistics.
   \item LLM Prompt Construction: Formats aggregate metrics into natural language prompt payload.
   \item Asynchronous API Call \& Caching: Communicates with Groq API (\texttt{llama-3.1-8b-instant}), parses response, and caches result (\texttt{CACHE\_DURATION\_MS = 2 minutes}).
 \end{itemize}

 \item \textbf{Phase 5: User Interface \& Verification}
 \begin{itemize}
   \item UI Implementation: HTML5, CSS3, Vanilla JavaScript (\texttt{fetch} API) interacting with Spring REST Controllers.
   \item Chart.js Rendering \& Verification: Interactive visualizations and end-to-end functional validation.
 \end{itemize}
\end{itemize}

\begin{figure}[ht]
\centering
\includegraphics[width=0.8\textwidth]{../Figures/system_architecture.png}
\caption{Finlytics System Architecture and Data Ingestion Flow}
\label{fig:system_architecture}
\end{figure}

\end{document}
